\documentclass[11pt,a4paper]{article}

\usepackage[margin=1in]{geometry}
\usepackage[hidelinks]{hyperref}
\usepackage{enumitem}

\setlength{\parindent}{0pt}
\setlist[itemize]{leftmargin=*, noitemsep, topsep=2pt}

\newcommand{\cvsection}[1]{%
  \vspace{10pt}
  {\large\textbf{#1}}\\[-2pt]
  \hrule
  \vspace{6pt}
}

\begin{document}

\begin{center}
  {\LARGE \textbf{Vyacheslav Kovalev}}\\
  \vspace{4pt}
  +7 904 442-51-20 \quad
  Telegram: \href{https://t.me/HunchHunter}{@HunchHunter} \quad
  \textbf{Full portfolio: \href{https://slavoch.github.io/cv/}{https://slavoch.github.io/cv/}}\\
  \textit{Additional project details, publications, and updates are available in the portfolio.}
\end{center}

\cvsection{Professional Summary}
RL Engineer and robotics researcher with about 4 years of experience in reinforcement learning, model predictive control, and sim-to-real deployment for legged and humanoid robots. Experienced across the full robotics development cycle: simulation, policy training, C++/ROS 2 deployment, hardware validation, and publication-quality research.

\cvsection{Work Experience}
\textbf{RL Engineer} \hfill June 2024 -- Present\\
\textit{Laboratory of Wave Processes and Controlling Systems, MIPT}
\begin{itemize}
  \item Developed and trained reinforcement learning policies for humanoid robots, including dynamic walking, soccer skills, balance recovery, and trajectory-following dance motions.
  \item Built and tuned robotics simulations in NVIDIA Isaac Sim, Isaac Gym, MuJoCo, and PyBullet, including dynamics, contacts, and sensor configuration.
  \item Deployed trained policies to real robots using high-performance C++ and ROS/ROS 2 nodes, topics, services, and hardware integration pipelines.
  \item Reduced the sim-to-real gap for custom in-house robots using domain randomization, test-stand validation, sensor calibration, and noise filtering.
  \item Developed a MuJoCo XLA (MJX) differentiable-simulation method for robot parameter identification using motion tracks only, without torque sensors.
  \item Achieved a 75\% reduction in turning deviation and a 46\% increase in traveled distance on a real biped robot compared with the baseline RL policy.
  \item Improved trajectory tracking accuracy by 1.88x on custom high-gear-ratio actuators, reducing MAE from 14.20 to 7.54 mrad.
\end{itemize}

\textbf{Research Assistant} \hfill August 2021 -- May 2023\\
\textit{Skolkovo Institute of Science and Technology}
\begin{itemize}
  \item Developed a hybrid MPC + RL control algorithm for a quadruped robot, serving as the main developer of the MPC component.
  \item Built the robot control dynamics model and implemented model predictive control in C++ with kinematic and dynamic constraints.
  \item Completed sim-to-real transfer from Gazebo simulation to physical hardware using RL policy training, domain randomization, deployment, and real-robot validation.
  \item Integrated control modules through ROS/ROS 2 nodes, topics, services, and actions for stable and safe operation on physical hardware.
\end{itemize}

\cvsection{Education}
\textbf{PhD Candidate, System Analysis, Control, Information Processing, and Statistics} \hfill Expected 2027\\
Moscow Institute of Physics and Technology

\vspace{4pt}
\textbf{MSc, Information Systems and Technologies} \hfill 2023\\
Skolkovo Institute of Science and Technology

\newpage

\cvsection{Skills}
\textbf{Advanced:} Python, PyTorch, Reinforcement Learning, Sim-to-Real, MPC, NVIDIA Isaac Sim, Isaac Gym, MuJoCo, Stable-Baselines3, PPO\\
\textbf{Intermediate:} Linux, C++, Git, ROS/ROS 2, PyBullet, Gazebo, Docker, skrl

\cvsection{Selected Publications}
\begin{itemize}
  \item \textbf{Achieving Precise and Reliable Locomotion with Differentiable Simulation-Based System Identification}. First-author IROS 2025 paper on real-robot locomotion and system identification: \url{https://ieeexplore.ieee.org/abstract/document/11247646}.
  \item \textbf{Trajectory-based actuator identification via differentiable simulation}. First-author IEEE Access paper on torque-sensor-free actuator identification: \url{https://arxiv.org/abs/2604.10351v1}.
  \item \textbf{Combining model-predictive control and predictive reinforcement learning for stable quadrupedal robot locomotion}. Preprint on hybrid MPC + RL control for quadruped robots: \url{https://arxiv.org/abs/2307.07752}.
\end{itemize}

\cvsection{Selected Project}
\textbf{Differentiable Simulation for Robot System Identification}\\
\href{https://wavegit.mipt.ru/Slavoch/mjx_sysid}{wavegit.mipt.ru/Slavoch/mjx\_sysid}
\begin{itemize}
  \item Developed a system identification approach for custom robotic hardware using differentiable simulation and real motion tracks.
\end{itemize}

\end{document}
